% Part: incompleteness
% Chapter: representability-in-q
% Section: representing-relations

\documentclass[../../../include/open-logic-section]{subfiles}

\begin{document}

\olfileid{inc}{req}{rel}

\olsection{సంబంధాలకు ప్రాతినిధ్యం చేయడం}

ఒక \emph{సంబంధం} ప్రాతినిధ్యం చేయదగినది అనడం అంటే ఏమిటో చెప్పుకుందాం.

\begin{defn}
\ollabel{defn:representing-relations} సహజ సంఖ్యలపై సంబంధం
$R(x_0,\dots,x_k)$ను కింది షరతు నెరవేరితే~$\Th{Q}$లో
\emph{ప్రాతినిధ్యం చేయదగినది} అంటాం: $R(n_0,\dots,n_k)$ సత్యమైనప్పుడల్లా
$\Th{Q}$ అనేది $!A_R(\num{n_0},\dots,\num{n_k})$ను, $R(n_0,\dots,n_k)$ అసత్యమైనప్పుడల్లా
$\Th{Q}$ అనేది $\lnot !A_R(\num{n_0},\dots,\num{n_k})$ను నిరూపించేలా ఒక
సూత్రం $!A_R(x_0,\dots,x_k)$ ఉండాలి.
\end{defn}

\begin{thm}
\ollabel{thm:representing-rels} ఒక సంబంధం~$\Th{Q}$లో ప్రాతినిధ్యం
చేయదగినదైతే, అప్పుడు మాత్రమే అది గణనీయమైనది.
\end{thm}

\begin{proof}
ముందుకు దిశ కోసం $R(x_0,\dots,x_k)$కు సూత్రం
$!A_R(x_0,\dots,x_k)$ ప్రాతినిధ్యం చేస్తుందని అనుకుందాం. $R$ను గణించే
క్రమవిధానం ఇది: $n_0$,~\dots,~$n_k$ నివేశంపై
$!A_R(\num{n_0},\dots,\num{n_k})$ నిరూపణ కోసం, దాని నిషేధం
$\lnot !A_R(\num{n_0},\dots,\num{n_k})$ నిరూపణ కోసం ఏకకాలంలో అన్వేషించండి.
మన పరికల్పన వల్ల అన్వేషణ వాటిలో ఒకదాన్ని తప్పక కనుగొంటుంది. మొదటిదైతే
``అవును'' అని, లేకపోతే ``కాదు'' అని నివేదించండి.

మరో దిశలో $R(x_0,\dots,x_k)$ గణనీయమైనదని అనుకుందాం. నిర్వచనం ప్రకారం దీని
అర్థం లక్షణ ప్రమేయం $\Char{R}(x_0,\dots,x_k)$ గణనీయమైనదని.
\olref[int]{thm:representable-iff-comp} వల్ల $\Char{R}$కు ఏదో సూత్రం
$!A_{\Char{R}}(x_0,\dots,x_k,y)$ ప్రాతినిధ్యం చేస్తుంది.
$!A_R(x_0,\dots,x_k)$ను $!A_{\Char{R}}(x_0,\dots,x_k,\num{1})$గా
తీసుకోండి. ఏ $n_0$,~\dots,~$n_k$లకైనా $R(n_0,\dots,n_k)$ సత్యమైతే
$\Char{R}(n_0,\dots,n_k)=1$; అప్పుడు $\Th{Q}$ అనేది
$!A_{\Char{R}}(\num{n_0},\dots,\num{n_k},\num{1})$ను, అందువల్ల
$!A_R(\num{n_0},\dots,\num{n_k})$ను నిరూపిస్తుంది. మరోవైపు
$R(n_0,\dots,n_k)$ అసత్యమైతే $\Char{R}(n_0,\dots,n_k)=0$. దీని అర్థం
$\Th{Q}$ కిందిదాన్ని నిరూపిస్తుంది:
\[
\lforall[y][(!A_{\Char{R}}(\num{n_0}, \dots, \num{n_k}, y) \lif y =
  \num{0})].
\]
$\Th{Q}$ అనేది $\eq/[\num{0}][\num{1}]$ను నిరూపిస్తుంది కనుక,
$\Th{Q}$ అనేది $\lnot !A_{\Char{R}}(\num{n_0},\dots,\num{n_k},\num{1})$ను,
అందువల్ల $\Th{Q}$ అనేది $\lnot !A_R(\num{n_0},\dots,\num{n_k})$ను కూడా నిరూపిస్తుంది.
\end{proof}

\begin{prob}
$R$ అనేది~$\Th{Q}$లో ప్రాతినిధ్యం చేయదగినదైతే, $\Char{R}$ కూడా అలాంటిదేనని
చూపండి.
\end{prob}

\end{document}
